\section{検証 I：$[1]$}\label{sec:one}

系~\ref{cor:one} により，ホモロジーを入力とする限り，任意の結び目・任意の framing で $[1]$ は正である．
この節では，それを結び目表で確かめ，$2J_r$ の偶奇の振る舞いを調べる．

% [R1]
\begin{compresult}\label{cr:table}
11 交点以下の全素結び目 801 個（thin で交代的なもの 562 個，thin で非交代的なもの 143 個，thick なもの 96 個）で，
$\hf_{[1]}$ の $G$ 展開は整数で閉じ，展開の前も後も非負である．
\end{compresult}

ケーブルの $[1]$ については，注意~\ref{rem:tilde} のとおり，計算した値が非負で $\sigma$ 不変である．したがって命題~\ref{prop:one} を直接適用できる．
一般の代数的な iterated torus 結び目については，Cherednik--Danilenko の予想（\S\ref{sec:conjB}）の下で同じ結論が従う．

\subsection{$2J_r$ の偶奇と parity knot}
KN は，トーラス結び目では $(\rho,g,\beta)$ を固定すると $2J_r$ の偶奇がそろうことを観察し，これを $S_R$ の自明な作用の証拠とした．
$[1]$ ではこの条件を入力の言葉で書ける．

% [F14]
\begin{proposition}\label{prop:parity}
命題~\ref{prop:one} の仮定の下で，$\hN_{[1],g,\beta,J_r}$ の $2J_r$ の偶奇が $(g,\beta)$ ごとにそろうことは，次と同値である．
$P$ の $\ba$ の次数が連続する範囲ごとに，$e_\bt+e_\ba/2$ の偶奇が一定である．
$\ba$ の次数の台に隙間が無いとき，これは Chandler--Gorsky~\cite{Chandler22} の parity knot の条件（HHH 全体で $t_{\mathrm{DGR}}-a/2$ の偶奇が一定）と同じである．
\end{proposition}
\begin{proof}
$\hf_{[1]}$ の $\ba^{2\beta}$ の行には $P$ の隣り合う二つの行だけが入る．
$p_k$ の $G^g$ の係数に現れる $\bt$ の指数は $k-g$ と同じ偶奇をもつので，$2J_r\equiv e_\bt-g$ である．
$\ba$ の次数は偶数なので，$\pm e_\ba/2$ の偶奇は同じである．
\end{proof}

\begin{compresult}
801 個の全てで $\ba$ の次数の台に隙間は無く，$2J_r$ の偶奇がそろう結び目は 9 個である．
$\mathbf{3_1},\mathbf{5_1},\mathbf{7_1},\mathbf{9_1},\mathbf{11a_{367}},\mathbf{8_{19}},\mathbf{10_{124}}$（トーラス結び目）と，$\mathbf{10_{125}}$（thick），$\mathbf{10_{140}}$（thin）である．
これらは Chandler--Gorsky の定義による parity knot と一致する．
\end{compresult}

したがって，$2J_r$ の偶奇の一致は $S_R$ の作用の自明さそのものではなく，ホモロジーの parity の性質を反映している．
非トーラス結び目にも偶奇がそろう例がある．
parity knot の条件そのものは Chandler--Gorsky による．BPS 不変量の偶奇との同値性と，11 交点以下の一覧は本論文で与える．
